\documentclass[11pt]{article}
\usepackage[margin=0.95in]{geometry}
\usepackage[expansion=false]{microtype}
\usepackage{hyperref}
\usepackage{xcolor}
\hypersetup{colorlinks=true, linkcolor=blue!60!black, citecolor=teal,
  urlcolor=blue!60!black}
\usepackage{amsmath}
\usepackage{enumitem}
\setlist{itemsep=1pt, topsep=2pt}

\title{\bfseries Memory as a Learned Policy\\
\large A research proposal for compounding-knowledge AI agents}
\author{Parth Kocheta \\\small University of Maryland \;\textbar\; \texttt{pranayko021@gmail.com}}
\date{Submitted to the Pace Fellowship, Summer 2026}

\begin{document}
\maketitle
\vspace{-1em}

\section*{The thesis}

The most important constraint on AI-for-science is not model intelligence.
It is memory. A research scientist running a six-month investigation
accumulates context across hundreds of papers, dozens of failed
experiments, and a steadily evolving model of which lab techniques work
in their hands. A frontier language model running the same investigation
forgets what it read on Monday by Wednesday. The compounding asset of a
scientific career---knowing more about your specific subfield each week
than you did the week before---has no operating system in current AI
agents.

Every shipping memory system today (Mem0, Letta, Mastra, Zep, my own
prior knowledge-graph system) pairs an external store with a hand-tuned
control layer that decides what to write and how to read. The store grows
from use; the control layer does not. I think this is the layer that
needs to be \emph{learned}.

\section*{The proposal}

I am training two LoRA adapters on a 4B-parameter open-source language
model that share an op vocabulary across read and write sides of memory.
The \emph{write policy} ingests one conversation turn at a time and emits
operations into a typed knowledge graph---creates, updates, marks of
contradiction, merges, archivals. The \emph{read policy} retrieves from
the same graph and answers questions. Both policies are trained with
GRPO~\cite{shao2024deepseekmath} (DeepSeek's critic-free RL algorithm)
on the LoCoMo~\cite{maharana2024locomo} and
LongMemEval~\cite{wu2024longmemeval} benchmarks of long-horizon
conversational memory.

The hard problem is the write reward. Outcome supervision is sparse and
delayed: a write op only matters if the read policy can later use what
was stored. My contribution is a dense companion signal I call
\emph{evidence coverage}---the fraction of dialogue turns required to
answer a held-out question that the policy actually preserved in the
graph, computed deterministically from existing benchmark annotations
and the graph's provenance fields. Coverage carries 60\% of the reward;
an LLM-judge-anchored answer-correctness term carries 40\%. The mix
trains $\sim$6$\times$ faster per step and is robust to judge noise.

The same trained adapter compiles to three storage backends: a flat
chunk store, a typed-transition knowledge graph, and a Mastra-style
observation log. Backend transfer (one policy, three stores) is the
empirical claim that lets this work be a \emph{layer} rather than a
\emph{system}.

\section*{Why this matters at the intersection Pace cares about}

\paragraph{Infrastructure.} Long-horizon agents are bottlenecked on
context window economics. A learned memory operating system pushes the
expensive recall work into a small adapter ($\sim$200 MB) and a
substrate-agnostic vocabulary, decoupling the agent's intelligence from
the linear-in-history cost of every retrieval. This is the missing layer
between today's RAG stacks and a real ``personal compute substrate.''

\paragraph{Economics.} The value of an AI collaborator scales with how
much of your specific context it has compounded. Today the practical
ceiling is the context window. Cracking the memory-OS problem changes
the unit economics of customer support, sales engineering, scientific
research, and personal AI: the marginal hour with a user gets cheaper
because the model accumulates instead of re-deriving. Whoever owns the
adapter gets a network effect that the foundation-model layer does not.

\paragraph{The physical world.} The benchmark-shaped reward I rely on
(LoCoMo's per-question evidence labels) is a research crutch. Real
deployments have no such labels. The fellowship period would let me
investigate two routes to closing that gap: (a) weak evidence mining
where an LLM extracts dia-id labels from a session transcript at
end-of-conversation, and (b) next-turn implicit reward where the user's
follow-up message scores the prior turn's stored facts. Either route
turns the technique into something that can run in a deployed system
where users do not annotate.

\section*{Status}

The full system is implemented and committed to a public repository. The
write/read environments, the dense reward, the three backends, and the
co-training outer loop are all production code. Smoke runs through
Tinker (Thinking Machines' managed RL training service) confirm
end-to-end training, with coverage-mean of 0.77 and judge accuracy of
0.55 on a tiny 5-step smoke. The full Phase B run is queued. A 15-page
draft of the long-form paper exists; this proposal is the compressed
version.

\section*{What I would do at Pace}

The fellowship period would convert two pieces of this work into
publishable artifacts: (1) the empirical paper with the full backend-
transfer table on LoCoMo and LongMemEval, and (2) a short essay arguing
that memory is the next vertical stack of LLM-agent infrastructure
worth investing in---with a clear taxonomy of what is durable
(substrate-agnostic adapters), what is interchangeable (the storage
itself), and what gets disrupted (per-customer fine-tuning).

\bigskip
\noindent\textit{All code, paper, and run logs are public:
\href{https://github.com/pkempire/pie}{github.com/pkempire/pie}.}

\bibliographystyle{plain}
\begin{thebibliography}{9}
\bibitem{shao2024deepseekmath} Z.\ Shao et al. ``DeepSeekMath: Pushing
the limits of mathematical reasoning in open language models.''
\emph{arXiv:2402.03300}, 2024.
\bibitem{maharana2024locomo} A.\ Maharana et al. ``Evaluating very
long-term conversational memory of LLM agents.''
\emph{arXiv:2402.17753}, 2024.
\bibitem{wu2024longmemeval} D.\ Wu et al. ``LongMemEval: Benchmarking
chat assistants on long-term interactive memory.''
\emph{arXiv:2410.10813}, 2024.
\end{thebibliography}

\end{document}
